\section{Limitations}
\label{sec:limitations}

\paragraph{Network representation.} The model has a single node. Internal transmission constraints, losses and
stability limits are absent. The four interconnections (Montenegro, Kosovo, North Macedonia, Greece) are
collapsed into one import and one export limit, so this is not a power-flow model. Regional correlation of
droughts and prices is represented only through the scenario availability factor $a_s$.

\paragraph{Market data.} Demand and prices are synthetic. Export revenue is a fixed 85\% of the import price.
Many conclusions, especially the value of PV, which in this model can always be exported up to 900~MW, depend
on this assumption. With observed ENTSO-E load and regional day-ahead prices the analysis could be repeated
without code changes.

\paragraph{Hydrology.} Inflows are reconstructed from reanalysis precipitation with an uncalibrated conceptual
model and scaled to an assumed mean energy. They are not gauged inflows. Drought is a uniform scaling of one
reference year. Real droughts also change the seasonal timing of inflows, and this is not captured.

\paragraph{Reservoir simplification.} The three-stage Drin cascade is one equivalent reservoir with a fixed
aggregate head. Head-dependent efficiency, travel times, individual reservoir limits and flood-control rules
are omitted. The other hydro plants are treated as run-of-river without storage.

\paragraph{Operational detail.} There is no unit commitment, no ramping, no reserve or ancillary-service
requirement and no thermal fleet. Batteries are therefore valued only for energy arbitrage and shortage
avoidance, which probably understates their real value. Simultaneous charging and discharging is excluded by
economics, not by binaries, and is monitored.

\paragraph{Time structure.} Investment and stochastic studies use eight representative weeks with
fleet-independent seasonal boundaries. Water allocation across seasons therefore does not adapt to new
investment. We quantify the error against full-year chronology (Section~\ref{sec:chrono}), but the planning
results inherit it.

\paragraph{Costs and uncertainty.} Technology costs are global IRENA averages with an assumed regional
multiplier. VOLL, the discount rate, lifetimes and the scenario probabilities are assumptions. Five scenarios
cannot describe the full climate distribution, and the 15\% stress probability is a judgement, not an estimate.

\paragraph{Validation.} The model has been \emph{verified}: against an independent PyPSA formulation and through
automated conservation and behaviour tests. It has \textbf{not been validated} against observed operations of
the Albanian system. Scenario comparisons are model experiments, not empirical evidence. The results are model
implications under stated assumptions and should not be read as policy recommendations.
